\documentclass[12pt]{article}
\usepackage[utf8]{inputenc}
\usepackage[T1]{fontenc}
\usepackage[margin=1in]{geometry}
\usepackage{setspace}
\usepackage{xcolor}
\usepackage[colorlinks=true, urlcolor=blue!60!black, linkcolor=blue!60!black]{hyperref}
\setlength{\parindent}{0em}
\setlength{\parskip}{0.6em}
\renewcommand{\baselinestretch}{1.0}
\begin{document}
\begin{center}
{\LARGE \textit{Snyder v. Phelps}, 562 U.S. 443 (2011)\par}
\vspace{0.5em}
{\large Case Brief --- Two-Page Edition\par}
\end{center}

\section{Memory Jogger}
Peaceful picketing on public concerns, on public land 1,000 feet from a Marine's funeral, cannot ground tort liability (\href{https://supreme.justia.com/cases/federal/us/562/443/}{562 U.S. 443}).

\section{Detailed Case Facts}
Westboro Baptist Church members, following police directions, picketed on public land about 1,000 feet from the Maryland funeral of Marine Lance Corporal Matthew Snyder, with signs such as ``Thank God for Dead Soldiers'' and ``God Hates Fags.'' His father saw only the tops of the signs.

\section{Procedural History}
A jury awarded Albert Snyder \$10.9 million, later reduced to \$5 million, for IIED, intrusion upon seclusion, and civil conspiracy. The Fourth Circuit reversed.

\section{Judicial Votes}
Affirmed 8--1. Roberts wrote for the Court; Breyer concurred; Alito dissented.

\section{Holding}
The First Amendment bars tort liability for Westboro's peaceful picketing on matters of public concern in a traditional public forum; the captive-audience doctrine does not apply.

\section{Analysis of Opinions}
\textbf{Majority:} Public-issue speech occupies the ``highest rung of the hierarchy of First Amendment values,'' judged by ``content, form, and context.'' The signs' ``overall thrust and dominant theme'' was public, Westboro had no prior grudge against the family, and the ``outrageousness'' standard invites juries to punish unpopular views. \textbf{Breyer:} the ruling is narrow and turns on proportionality. \textbf{Alito:} the First Amendment is no license for a ``vicious verbal assault'' on grieving private figures.

\section{Examples: Future Applications}
\textbf{Same-Side:} activists who protest police violence from a distant sidewalk near an officer's funeral are protected. \textbf{Opposite-Side:} a content-neutral funeral buffer law likely survives. \textbf{Fence-Sitter:} picketing a grieving family's home and posting their address online.

\section{Critique}
Alito's paradox: targeting a stranger for publicity earns more protection than a private feud. And the ``dominant theme'' test may let a few political slogans immunize harassment.

\section{Key Quotations}
\begin{quote}
``Speech is powerful. It can stir people to action, move them to tears of both joy and sorrow, and---as it did here---inflict great pain. On the facts before us, we cannot react to that pain by punishing the speaker. As a Nation we have chosen a different course---to protect even hurtful speech on public issues to ensure that we do not stifle public debate.''
\end{quote}
\noindent --- Chief Justice Roberts, Opinion of the Court, 562 U.S. at 460--461.

\end{document}
